\section{Introduction and Contribution Statement}
\label{sec:prof-introduction}

Physical AI systems do not act on truth. They act on telemetry, inferred state, delayed
measurements, missing packets, stale buffers, and learned surrogates of the world. The
resulting engineering problem is deeper than forecast error or planner suboptimality:
the software stack can remain internally coherent while the physical system becomes
externally unsafe. Batteries may violate hidden state-of-charge limits, vehicles may
enter unsafe headway regimes on stale perception, industrial plants may cross thermal
or pressure envelopes, clinical monitoring pipelines may suppress timely intervention,
and navigation or flight systems may remain locally consistent while drifting outside a
feasible corridor.

ORIUS is built around the claim that these failures share one universal object:
degraded observation changes the semantics of what may still be safely done. The field
therefore needs a runtime layer whose job is not to outperform every planner, but to
mediate the boundary between degraded observation and physical actuation. ORIUS,
Observation--Reality Integrity for Universal Safety, is proposed as that layer. It
detects reliability loss, calibrates the observation-consistent state set, constrains
the admissible action set, repairs or replaces unsafe candidate actions, and certifies
release through an auditable lifecycle. The framework is written around the universal
runtime claim rather than a witness-row-specific framing. Battery is retained as the
deepest witness row, not as the conceptual center of the claim.

This manuscript positions ORIUS against a precise gap in the literature. Distribution-
free inference can quantify uncertainty
\cite{vovk2005algorithmic,lei2018distribution,romano2019conformalized,gibbs2021adaptive,angelopoulos2023gentle,barber2023beyond}.
Runtime assurance and shields can intercept unsafe behavior
\cite{sha1998simplex,sha2001using,desai2019soter,bloem2015shield,bartocci2018introduction,leucker2009brief}.
Safety filters and constrained control can map uncertainty into admissible action sets
\cite{ames2019control,ames2014cbf,wabersich2021predictive,fisac2019general,mayne2005robust,rawlings2017mpc}.
Fault, anomaly, and drift detection can identify when observation should no longer be
trusted \cite{basseville1993detection,page1954continuous,page1955test,hinkley1971inference,pasqualetti2013attack,teixeira2015secure}.
What none of these families alone provides is a reusable runtime grammar that links
degraded observation to repaired actuation and auditable certificate release across
multiple physical domains.

The paper therefore makes four contributions.
\begin{enumerate}[leftmargin=1.25em,itemsep=0.2em]
  \item It defines the observation--action safety gap (OASG) as the right universal
  hazard object for Physical AI under degraded observation.
  \item It presents the Detect--Calibrate--Constrain--Shield--Certify kernel as a typed
  runtime safety layer with explicit formulas, pseudocode, and certificate semantics.
  \item It unifies domain evaluation through one adapter contract, one benchmark schema,
  and one CertOS governance surface.
  \item It argues for a bounded-universal evidence posture: battery is the witness row;
  AV, industrial, and healthcare are defended bounded rows; navigation and aerospace
  remain explicitly gated.
\end{enumerate}

The draft is intentionally ambitious but disciplined. It uses a universal writing
posture, but it keeps one explicit evidence-boundary section in the main body. It
includes formulas, theorem statements, proof sketches, runtime pseudocode, decisive
results tables, and cross-domain synthesis, while pushing proof-depth and artifact-depth
material into two separate appendices.
